\section{The exact logical substrate}

\subsection{Ordinary iPRA as the canonical interface}

The canonical target is the ordinary deeply embedded judgment
\[
  \Gamma\ \vdash_{\PRAAx}\ \varphi,
\]
implemented by \code{LeanProofs.ConciseIPRA.Prf}.  A sequent lives over a fixed finite number \(n\) of free numerical variables.  Its terms have type \(\PTm\,n\), its formulas have type \(\Fm\,\PTm\,n\), and every formula in \(\Gamma\) has the same scope.

Primitive-recursive function symbols are generated by zero, successor, projections, composition, and primitive recursion.  The term language has one application constructor, so every fixed primitive-recursive algorithm can be represented by a symbol in \(\PR\) and applied to arbitrary object terms.  The axiom family \(\PRAAx\) consists of defining equations for those symbols, successor separation and injectivity, and induction for quantifier-free formulas.

The logical calculus is intuitionistic natural deduction with equality.  The structural operations used by SCOPE are already checked:

\begin{itemize}
\item \code{Prf.subst}: simultaneous substitution of arbitrary object terms for free variables;
\item \code{Prf.weaken}: context weakening;
\item \code{Prf.trans}: simultaneous cut, replacing every assumption by a derivation of it;
\item \code{Prf.exE}: existential elimination with the witness kept under a fresh-variable scope;
\item \code{Prf.eqSubst}: simultaneous Leibniz transport.
\end{itemize}

These operations are sufficient for the entire structural part of SCOPE.  No new logical axiom, choice principle, or reflection rule is needed.

\subsection{Quantifier-free induction and the compilation discipline}

The axiom \code{PRA.qfInd} applies to a formula \(I(t,\vec x)\) only after Lean has established \code{I.QF}.  The checked derivation library exposes a convenient open form \code{LeanProofs.IPRA.ind}.  SCOPE uses this restriction as a design rule rather than treating it as an inconvenience:

\begin{principle}[Equation-first induction]
Any finite property needed inside an induction is compiled to a primitive-recursive bit or state component.  The induction matrix is an equation or a finite quantifier-free combination of equations.  Reflection back to the intended bounded formula happens before or after the induction, never by pretending a bounded quantifier is syntactically quantifier-free.
\end{principle}

For example, ``every table entry below \(t\) is valid'' is represented by the equation \(\mathsf{all}(t,q)=1\).  ``Some transition has failed'' is a failure-accumulator bit.  ``Every coordinate is within its cap'' is either a fixed finite conjunction or a finite-fold bit.  This uniform pattern is what makes trace compilation reusable.

\subsection{Current arithmetic support}

The present \code{IPRADerivations} library is substantially richer than several earlier framework packets assumed.  SCOPE relies on the following checked components.

\begin{itemize}
\item \code{Core.Basic}: arithmetic symbols, defining equations, open induction, and instantiation;
\item \code{Core.Definitions}: \code{PTm.toPR}, its evaluation theorem \code{toPR_eq}, arbitrary-arity \code{recDef}, and finite sums/products;
\item \code{Core.Order}: reflected order, totality, equality decidability, and monotonicity;
\item \code{Core.Pairing}: \code{eqb}, pairing/projections, and simultaneous recursion;
\item \code{Core.Search}: a least-zero search \(\mu\), soundness, foundness, and minimality;
\item \code{Core.Division}, \code{Core.MultiSums}, and \code{Core.Permutation}: the arithmetic needed for finite code spaces;
\item polynomial and integer-difference backends, finite descent, and finite-dimensional linear algebra for application nodes.
\end{itemize}

One reported obstacle is therefore obsolete.  Older bounded-reflection outlines asked for an arity-indexed abstraction turning an open term into a primitive-recursive symbol.  The current declaration
\[
  \code{PTm.toPR} : \PTm\,k \longrightarrow \PR\,k
\]
and its object-theory theorem \code{toPR_eq} provide exactly that operation.  Together with \code{recDef}, they permit arbitrary-parameter Boolean folds without an internal universal evaluator.

\subsection{Fixed specialization, not a universal evaluator}

SCOPE is a schema over a \emph{fixed} finite program or formula.  Lean's metatheory traverses the finite syntax tree, constructs a particular \(\PR\) symbol, and constructs a particular \(\Prf\,\PRAAx\) derivation.  The object theory need not evaluate arbitrary encoded \(\PR\) programs.

This distinction avoids a common impossibility.  A universal total evaluator for all primitive-recursive programs cannot itself be primitive recursive.  Nothing in SCOPE asks for one.  Each theorem has finitely many concrete symbols, and \code{PTm.toPR} compiles those concrete open terms directly.

\subsection{Object proofs versus semantic Lean theorems}

A native Lean theorem about \(\Nat\), lists, rationals, or a recursive function is useful for discovering an invariant and testing an algorithm.  It is not automatically a theorem of the embedded object theory.  Soundness goes from an object derivation to semantics; there is no unrestricted reverse reflection from arbitrary Lean proofs.

Accordingly, every SCOPE node has two logically distinct ingredients:

\begin{enumerate}
\item a host-level finite description used to generate symbols, terms, and scheduling data;
\item an object-level derivation builder whose output type is a \(\Prf\,\PRAAx\) sequent.
\end{enumerate}

The compiler is trusted only as a producer of a proof term that Lean's kernel checks.  External code may propose certificates and can greatly accelerate discovery, but acceptance is ultimately represented by object arithmetic and an object derivation.

\subsection{Backend transport}

The repository's bridge now exposes closed, open, and scoped transports between ordinary ConciseIPRA and Foundation's standard iPRA presentation.  The relevant APIs include \code{standardToConciseSequent}, \code{conciseToStandardSequent}, and scoped round-trip theorems.  Thus open-sequent transport is no longer a missing generic theorem.

SCOPE therefore develops all reusable mathematics once in the ordinary \(\Prf\,\PRAAx\) interface.  Native scalar calculi and Foundation arithmetic are export backends.  This avoids parallel libraries and lets application endpoints choose the presentation they require.
